\chapter{Meshing with gmsh}\label{ch:meshing}

Most numerical solvers do not work on NURBS geometry directly. Finite-element (FEM),
boundary-element (BEM) and finite-volume codes need a mesh of triangles, quadrangles or
tetrahedra. \NF{} does not mesh by itself. It hands its geometry \emph{exactly} to
\href{https://gmsh.info}{gmsh}, an established open-source mesh generator. The subpackage
\code{nurbforge.mesh} contains this bridge. This chapter covers:
\begin{itemize}
  \item the exact transfer of curves, surfaces and trimmed surfaces;
  \item shells and solids, with physical groups taken from face labels;
  \item local mesh sizes and three-dimensional meshes;
  \item BEM-ready surface meshes as NumPy arrays;
  \item file output, and the implementation details that matter in practice.
\end{itemize}

\section{Installation and the gmsh API}

gmsh is an optional dependency:

\begin{lstlisting}[style=shell]
pip install "nurbforge[mesh]"     # or: pip install gmsh  (version 4.11 or newer)
\end{lstlisting}

Without gmsh, importing \code{nurbforge.mesh} still works. The first function that needs gmsh
raises an \code{ImportError} that explains how to install it. On AlmaLinux and RHEL the pip wheel of
gmsh needs a few X11 and OpenGL runtime libraries, even when it is used without a display (see
\cref{ch:install}). If they are missing, \code{import gmsh} raises \code{OSError}, not
\code{ImportError}. Scripts that treat gmsh as optional therefore catch both:

\begin{lstlisting}[style=python]
try:
    mesh_surfaces([front, back], "lens.msh", size=1.0)
except (ImportError, OSError) as exc:    # gmsh missing or its libraries unavailable
    print(f"meshing skipped: {exc}")
\end{lstlisting}

gmsh is a stateful library: there is one global model, which is created with
\code{gmsh.initialize()} and released with \code{gmsh.finalize()}. \NF{}'s functions work on
this current model. The context manager \apiref{mesh.GmshSession} handles initialisation and
clean-up. Inside the session the complete gmsh Python API is available as \code{g.gmsh} (the
imported \code{gmsh} module itself), so anything not wrapped by \NF{} can be done with gmsh
calls directly.

\section{Exact transfer of NURBS geometry}\label{sec:mesh-transfer}

\subsection{What ``exact'' means}

\NF{} passes the original NURBS data to gmsh's OpenCASCADE (OCC) kernel, through
\code{occ.addBSpline} and \code{occ.addBSplineSurface}: the degree, the distinct knots with
their multiplicities, the control points and the weights. gmsh therefore meshes the true
geometry, not a tessellation. Mesh nodes lie on the exact surface, so a sphere stays a sphere,
whatever the mesh size. The test suite checks that gmsh's own evaluation
(\code{gmsh.model.getValue}) agrees with \NF{} to $10^{-9}$. In the example below the
difference is $5\cdot10^{-16}$, which is rounding.

Two technical details are handled automatically:
\begin{itemize}
  \item \textbf{Control point order.} \NF{} stores a surface's control net row-major, with
    $v$ running fastest (index $i\cdot n_v+j$). gmsh expects $u$ running fastest.
    \code{add_surface} reorders the points.
  \item \textbf{Temporary points.} The control points are created as gmsh points and removed
    again after the spline has been built. Otherwise they would stay in the model as free
    vertices and end up as stray mesh nodes.
\end{itemize}

\subsection{The transfer functions}

\begin{center}
\begin{tabularx}{\textwidth}{@{}lX@{}}
\toprule
Function & Effect \\
\midrule
\apiref{mesh.add_curve}\code{(curve)} & Adds a \code{NurbsCurve} as a B-spline curve and
  returns the gmsh curve tag. Useful for 1D meshes and embedded lines. \\
\apiref{mesh.add_surface}\code{(surface, wire_tags=None)} & Adds a \code{NurbsSurface} and
  returns the surface tag. \code{wire_tags} optionally trims the surface with gmsh wires
  defined in its \emph{parameter} space, outer wire first. \\
\apiref{mesh.add_trimmed_surface}\code{(trimmed)} & Adds a \code{TrimmedSurface}. The trim
  loops are passed exactly, as B-spline wires in parameter space, outer loop first, then the
  holes. The helper curves are removed afterwards. An untrimmed face is passed as the plain
  surface, so OCC can handle a closed (periodic) patch itself. \\
\apiref{mesh.add_shell}\code{(shell, ...)} & Adds a shell with its topology and, optionally,
  a volume and physical groups (\cref{sec:mesh-shell}). \\
\bottomrule
\end{tabularx}
\end{center}

These module functions work on whatever gmsh model is current. The session methods
\apiref{mesh.GmshSession.add_curve} and \apiref{mesh.GmshSession.add_surface} do the same, and
\code{add_surface} dispatches on the type: a \code{NurbsSurface} goes to \code{add_surface},
and a \code{TrimmedSurface} to \code{add_trimmed_surface}.

The following script transfers a sphere (a single rational surface), a plate with a circular
hole (a trimmed surface) and a circle. It compares gmsh's evaluation with \NF{}'s, meshes the
surfaces and writes the mesh in two formats.

\example{15_mesh_exact_transfer.py}

\begin{lstlisting}[style=text]
max |gmsh - nurbforge| on the sphere: 4.58e-16
sphere: 2110 triangles
plate: 159 triangles
sphere nodes: radius between 2.000000000000 and 2.000000000000
wrote transfer.msh and transfer.vtk
\end{lstlisting}

Every node of the sphere mesh lies at radius $2$ to twelve decimal places. For trimmed
surfaces, nodes on the trim edges follow OCC's own 3D approximation of the trim curve. Their
deviation is about $10^{-4}$ in the tests. Interior nodes are exact.

\begin{note}[Separate surfaces do not share edges]
Surfaces added one by one are independent in gmsh. Even if two of them touch along a common
boundary, each has its own boundary curve. \code{mesh(heal=True)} calls
\code{occ.removeAllDuplicates()}, which merges coincident edges of simple adjacent patches. For
real B-rep models, always use \code{add_shell}. It keeps the topology, including seams, as
described in the next section.
\end{note}

\section{The session: \texttt{GmshSession}}\label{sec:mesh-session}

\begin{lstlisting}[style=python]
from nurbforge.mesh import GmshSession

with GmshSession(model_name="nurbforge", verbose=False) as g:
    g.add_shell(part, physical_groups=True)
    g.refine("inlet", 0.05, distance=1.0)
    g.mesh(dim=3, size=0.3, order=2, curvature=24)
    g.write("part.msh")
\end{lstlisting}

\begin{longtable}{@{}>{\raggedright\arraybackslash}p{0.3\textwidth}p{0.64\textwidth}@{}}
\toprule
\textbf{Member} & \textbf{Description} \\
\midrule\endhead
\code{GmshSession(model_name, verbose)} & Constructor (defaults \code{"nurbforge"} and \code{False}). On entering the
  \code{with} block, gmsh is initialised (unless it already is) and a model of this name is
  added. \code{verbose=True} turns on gmsh's terminal output. Leaving the block always calls
  \code{gmsh.finalize()}, also after an exception. \\
\code{g.gmsh} & The gmsh module, for direct API calls. \\
\code{add_curve(curve)} & See \cref{sec:mesh-transfer}. \\
\code{add_surface(surface)} & \code{NurbsSurface} or \code{TrimmedSurface}. Returns the tag. \\
\code{add_shell(shell, solid=True, physical_groups=False, volume_name="volume")} & Adds a shell
  (\cref{sec:mesh-shell}). Returns the surface tags in face order, without repetitions,
  followed by the volume tag if a volume was created. The session also remembers the face
  labels. \\
\code{face_tags()} & Dictionary mapping each face label to its current gmsh surface tags.
  Unlabelled faces are listed under \code{"unlabeled"}. \\
\code{refine(where, size, distance=0.0)} & Requests a local mesh size on faces
  (\cref{sec:mesh-sizes}). \\
\code{mesh(dim=2, size=None, order=1, heal=True, curvature=None)} & Generates the mesh, see
  below. \\
\code{surface_mesh(orient)} & Returns the triangle surface mesh as a
  \code{SurfaceMesh}, checked and, with \code{orient=True} (default), oriented (\cref{sec:mesh-bem}). \\
\code{write(path)} & Writes the mesh or geometry. The format is chosen by the extension
  (\cref{sec:mesh-output}). \\
\bottomrule
\end{longtable}

\subsection{\texttt{mesh}: generating the mesh}

\begin{longtable}{@{}p{0.17\textwidth}p{0.12\textwidth}p{0.63\textwidth}@{}}
\toprule
\textbf{Parameter} & \textbf{Default} & \textbf{Meaning} \\
\midrule\endhead
\code{dim} & 2 & 1 meshes curves only, 2 adds triangles on the surfaces, 3 adds tetrahedra in
  the volumes. A volume exists only for closed shells added with \code{solid=True}. \\
\code{size} & \code{None} & Global maximum element size (\code{Mesh.MeshSizeMax}). Without
  it, gmsh chooses sizes from the model dimensions. \\
\code{order} & 1 & Element order. 2 gives quadratic elements (6-node triangles, 10-node
  tetrahedra). Higher orders are possible. The additional nodes are placed on the exact
  geometry. \\
\code{heal} & \code{True} & Calls \code{occ.removeAllDuplicates()} before meshing. This merges
  coincident points and edges of separately added surfaces. Shells added with
  \code{add_shell} already have a consistent topology. \\
\code{curvature} & \code{None} & Number of elements per $2\pi$ of curvature
  (\code{Mesh.MeshSizeFromCurvature}). Strongly curved regions, such as small fillets and
  bores, get smaller elements, while flat regions keep the global size. Must be at least 1. \\
\bottomrule
\end{longtable}

\code{mesh} first synchronises the OCC model, then sets the options and rebuilds the size
fields from all \code{refine} calls. Finally it calls \code{gmsh.model.mesh.generate(dim)},
and \code{setOrder(order)} if \code{order} $>1$. Calling \code{mesh} again re-meshes the
model with the current settings.

\section{Shells, solids and physical groups}\label{sec:mesh-shell}

\subsection{\texttt{add\_shell}}

\code{add_shell(shell, solid=True, physical_groups=False, volume_name="volume")} transfers a
complete B-rep shell. Its faces share their edges in gmsh as they do in \NF{}, so the mesh is
conforming: neighbouring faces have the same nodes along their common edge. For a
\emph{closed} shell with \code{solid=True}, a gmsh volume is created as well, which
\code{mesh(dim=3)} fills with tetrahedra. With \code{solid=False} only the surfaces remain.
That is what surface meshes for BEM need, because a volume would only be in the way. An
open shell never gets a volume.

The return value is the list of gmsh surface tags in the order of \code{shell.faces}, without
repetitions, followed by the volume tag if one was created.

\subsection{Physical groups from face labels}

Solvers assign boundary conditions and materials through \emph{physical groups}: named sets of
mesh entities that are written into the mesh file. \NF{} creates them from face labels.

Every \code{TrimmedSurface} carries a \code{label}:
\begin{itemize}
  \item The solid primitives label their faces: \code{bottom}, \code{top}, \code{front},
    \code{back}, \code{left} and \code{right} for a box; \code{side}, \code{bottom} and
    \code{top} for a cylinder and a cone; \code{surface} for a sphere and a torus.
  \item \apiref{Shell.with_face_labels} sets the labels of all faces at once.
  \item Booleans, transforms and STEP keep labels. Every piece of a split face inherits the
    label of the face it came from.
\end{itemize}
A label given before a chain of operations therefore still identifies the inlet of a pipe
after a slot has been cut into it.

With \code{physical_groups=True}, \code{add_shell} creates:
\begin{itemize}
  \item one physical \emph{surface} group per label, named after the label and containing all
    faces with that label;
  \item a physical \emph{volume} group named \code{volume_name}, if a volume was created.
\end{itemize}

\code{face_tags()} returns the current gmsh tags of every label. Faces are matched to gmsh
surfaces by an interior point, not by tag number. The result is therefore correct even after
gmsh has renumbered surfaces, for example after \code{removeAllDuplicates}, or after further
shells were added. Only the session method \code{g.add_shell} records the labels. The module
function \code{add_shell} creates physical groups as well, but \code{face_tags} and
\code{refine} cannot see its labels.

\section{Local mesh sizes}\label{sec:mesh-sizes}

\code{refine(where, size, distance=0.0)} requests a mesh size of \code{size} on selected faces:

\begin{itemize}
  \item \code{where} is a face label, a gmsh surface tag, or a list mixing both. Labels refer
    to shells added with \code{g.add_shell}.
  \item With \code{distance} $=0$, only the faces themselves get the size. gmsh's
    \code{Constant} field applies it exactly on the faces and their boundaries.
  \item With \code{distance} $>0$, the size grows linearly from \code{size} on the faces to the
    global size (\code{mesh(size=...)}) at that distance, also into the volume. This uses a
    \code{Distance} field and a \code{Threshold} field. Without a global size, the target at
    \code{distance} is the diagonal of the model's bounding box.
  \item Several refinements combine by taking the smallest size everywhere (a \code{Min}
    field, which becomes the background mesh).
  \item Faces next to a refined face also become finer near the shared edge, because gmsh
    grades the mesh away from a fine boundary.
\end{itemize}

Refinements are only stored when \code{refine} is called. They are resolved to gmsh surfaces
and applied in \code{mesh()}. The order of \code{refine} and \code{add_shell} calls therefore
does not matter, and renumbering in between is harmless. Errors are reported early or clearly:
\begin{itemize}
  \item a non-positive \code{size}, a negative \code{distance} or an empty list raises
    \code{ValueError} immediately;
  \item an unknown label raises \code{KeyError} in \code{mesh()}, listing the available labels.
\end{itemize}

\section{A complete solid example}

The script below labels the faces of a pipe, cuts a slot with a Boolean and transfers the
result with physical groups. It refines the inlet with growth into the volume and the outlet
on the face only. It then meshes the volume with tetrahedra, using curvature-based sizing, and
compares the volume of the tetrahedral mesh with the exact volume.

\example{15_mesh_solid.py}

\begin{lstlisting}[style=text]
face labels: ['back', 'bottom', 'front', 'inlet', 'outlet', 'top', 'wall']
surfaces [1, 2, 3, 4, 5, 6, 7], volume 1
face_tags: {'wall': [1], 'inlet': [2], 'outlet': [3], 'bottom': [4], 'top': [5],
            'front': [6], 'back': [7]}
  physical group dim=2 'wall'
  ...
  physical group dim=2 'back'
  physical group dim=3 'fluid'
38652 tetrahedra, mesh volume 14.8665, exact 14.9133
wrote pipe.msh
\end{lstlisting}

The slot's faces carry the labels of the box that cut it (\code{bottom}, \code{top},
\code{front}, \code{back}). Its left and right faces lie outside the pipe and have vanished.
The linear tetrahedra cut the curved wall by chords, so the mesh volume is about 0.3\,\%
smaller than the exact volume. A finer mesh, or \code{order=2}, reduces the difference.

\section{BEM-ready surface meshes}\label{sec:mesh-bem}

Boundary-element solvers discretise only the surface of a body. They expect a mesh that is:
\begin{enumerate}
  \item \textbf{closed} (watertight): every edge is shared by exactly two triangles;
  \item \textbf{consistently oriented}: neighbouring triangles traverse their common edge in
    opposite directions, so the normals do not flip from one triangle to the next;
  \item \textbf{oriented outwards}: the normals point out of the material.
\end{enumerate}
A small orientation error, such as a single reversed face, silently changes the sign of
integrals over that face and spoils the solution. \NF{} therefore checks all three properties
and repairs the orientation.

\subsection{\texttt{bem\_mesh}: one call}

\begin{lstlisting}[style=python]
from nurbforge.mesh import bem_mesh

m = bem_mesh(shell, size=0.2, order=2, curvature=None, refine={"top": 0.1},
             path="part.msh")
\end{lstlisting}

\apiref{mesh.bem_mesh} opens a session, adds the shell with \code{solid=False}, applies the
refinements (a dictionary from label or surface tag to local size), meshes in 2D, extracts and
orients the surface mesh, and optionally writes it to \code{path}. If a path is given, physical
groups per face label are written as well. The result is a \apiref{mesh.SurfaceMesh}.

\subsection{\texttt{SurfaceMesh}}

\begin{longtable}{@{}>{\raggedright\arraybackslash}p{0.33\textwidth}p{0.61\textwidth}@{}}
\toprule
\textbf{Member} & \textbf{Content} \\
\midrule\endhead
\code{nodes} & \code{(N, 3)} float array of node coordinates. Only nodes used by triangles are
  included. \\
\code{triangles} & \code{(M, 3)} int array for linear elements, or \code{(M, 6)} for quadratic
  elements. Indices are 0-based rows of \code{nodes}. The node order is gmsh's: three corners,
  then the midpoints of the edges 0--1, 1--2 and 2--0. \\
\code{surfaces} & \code{(M,)} gmsh surface tag of every triangle. \\
\code{node_tags} & \code{(N,)} gmsh node tag of every row of \code{nodes}, to match the arrays
  with a written \code{.msh} file. \\
\code{groups} & Dictionary mapping each face label to the array of triangle indices with that
  label. Every triangle is in exactly one group. \\
\code{order} & 1 or 2. \\
\code{normals()} & \code{(M, 3)} unit normals of the flat corner triangles. \\
\code{areas()} & \code{(M,)} areas of the corner triangles. \\
\code{centroids()} & \code{(M, 3)} centroids of the corner triangles. \\
\code{quality()} & \code{(M,)} shape quality of the corner triangles: 1 equilateral, 0
  degenerate (\cref{sec:mesh-shrinkwrap}). \\
\code{write(path)} & Writes \code{.msh} (MSH 2.2 ASCII) or binary \code{.stl} without gmsh. \\
\code{volume()} & Enclosed volume by the divergence theorem. It is positive for an
  outward-oriented mesh, and cavities count negative. For quadratic meshes it uses the corner
  triangles, so it converges like a linear mesh. \\
\code{boundary_edges()} & \code{(K, 2)} corner-node edges used by only one triangle. Empty for a
  watertight mesh. \\
\code{is_watertight()} & \code{True} if every edge is used by exactly two triangles. \\
\code{is_consistently_oriented()} & \code{True} if every directed edge occurs once. \\
\bottomrule
\end{longtable}

\subsection{\texttt{surface\_mesh} and the orientation rules}

\code{g.surface_mesh(orient=True)} does the work behind \code{bem_mesh} in an existing session,
after \code{g.mesh(dim=2)}. It proceeds in four steps.

\begin{enumerate}
  \item \textbf{Collect.} The triangles of all meshed surfaces are gathered into one array and
    their nodes renumbered.
  \item \textbf{Check.} The function raises \apiref{mesh.MeshTopologyError}, a subclass of
    \code{ValueError}, in these cases:
    \begin{itemize}
      \item an edge is used by only one triangle. The shell is open, or surfaces added
        separately with \code{add_surface} do not share their edges.
      \item an edge is used by more than two triangles (non-manifold).
      \item the model contains quadrangles (\code{Mesh.RecombineAll}), mixed element orders,
        or no surface mesh at all.
      \item the arrangement of surfaces is not orientable.
    \end{itemize}
  \item \textbf{Make consistent.} gmsh orients every surface by its own parametrisation, and
    neighbouring surfaces may disagree. A breadth-first search over the shared edges decides,
    surface by surface, whether to flip it. After this step every shared edge is traversed in
    opposite directions.
  \item \textbf{Orient outwards.} Every closed connected component is turned so that its
    normals point \emph{out of the solid}. On the outer boundary they point away from the
    material. On a cavity they point into the void, which is also away from the material. The
    decision uses the signed volume of the component and how many other components enclose
    it, determined by a ray-parity test. A component inside an even number of others must have
    positive volume, one inside an odd number negative volume.
\end{enumerate}

The flips are also applied to the gmsh model, with \code{gmsh.model.mesh.reverse}. A
subsequent \code{g.write("part.msh")}, \code{.stl} or \code{.vtk} therefore contains the same,
correctly oriented mesh. With \code{orient=False}, step 2 is still performed, but gmsh's
orientation is returned unchanged. That is useful for diagnostics.

\begin{note}
Shells produced by \NF{} are already oriented outwards, and the transfer through STEP
(\cref{sec:mesh-impl}) preserves the orientation. For such shells the repair in steps 3 and 4
usually changes nothing. It is a safety net, for example for models built from separately
created faces, or when gmsh's mesh has been modified by hand.
\end{note}

\subsection{Example: a body with a cavity}

\example{15_mesh_bem.py}

\begin{lstlisting}[style=text]
nodes (8272, 3), triangles (4134, 6), order 2
groups: {'bottom': 460, 'top': 462, 'front': 462, 'back': 458, 'left': 462,
         'right': 458, 'surface': 1372}
watertight True, consistent True
mesh volume 59.8454, exact 59.8112
outer faces pointing outwards: True
cavity faces pointing to the centre: True
total area 108.5098, exact 108.5664
before: False
after: True volume 6.2530
MeshTopologyError: surface mesh is not watertight: 24 open edge(s) (used by one
triangle). Mesh a closed shell added with add_shell, so that its faces share their edges.
\end{lstlisting}

The cavity, the sphere labelled \code{surface}, is refined to half the global size. Its
normals point to the centre of the sphere, which is into the void, while the box's normals
point outwards. In the second part, one face of a cylinder is reversed on purpose with a raw
gmsh call. \code{surface_mesh()} detects and repairs this, and the repaired orientation is
written to \code{cylinder.stl}.

\begin{tip}[Passing the mesh to a solver]
Most BEM codes need exactly \code{nodes} and \code{triangles}, possibly with a region number
per triangle. A region array follows directly from \code{groups}:
\begin{lstlisting}[style=python]
region = np.empty(len(m.triangles), dtype=int)
for k, (label, idx) in enumerate(sorted(m.groups.items())):
    region[idx] = k
\end{lstlisting}
For quadratic BEM elements, check the node order your solver expects. gmsh's order (corners,
then edge midpoints 0--1, 1--2, 2--0) is the most common convention.
\end{tip}

\section{Shrink wrap: a closed hull around anything}\label{sec:mesh-shrinkwrap}

\apiref{mesh.shrink_wrap} builds one closed triangle mesh around shells and loose surfaces that
overlap, leave gaps or are open. It is a fallback where an exact Boolean or sewing is not
possible, a repair for imported geometry, and a way to get a simplified outer boundary. Unlike
everything else in Nurbforge it is \emph{approximate}: the result is a mesh accurate to about
the resolution, not exact geometry. It needs numpy only, not gmsh.

\begin{lstlisting}[style=python]
from nurbforge.mesh import shrink_wrap

hull = shrink_wrap([body_a, body_b, loose_surface],
                   resolution=0.2,     # grid cell: about the edge length and the accuracy
                   offset=0.5,         # distance of the hull from the input (default 0)
                   fill_holes=True)
hull.is_watertight(), hull.volume()    # a SurfaceMesh, normals pointing outwards
hull.write("envelope.msh")             # or .stl
\end{lstlisting}

The objects are \code{Shell}, \code{TrimmedSurface} or \code{NurbsSurface} objects, one or a
list. The result is a \apiref{mesh.SurfaceMesh} with linear triangles, watertight and oriented
outwards; \apiref{mesh.SurfaceMesh.write} writes it as \code{.msh} (MSH 2.2 ASCII) or binary
\code{.stl}.

\begin{itemize}
  \item Gaps and openings narrower than about $1.5\,\times$ \code{resolution} (plus twice the
    \code{offset}) are closed. Wider ones are not, and the hull then follows the inside of an
    open part as well. Parts further apart give separate closed components.
  \item With \code{offset=0} the hull lies on the input: on a sphere the nodes are within about
    $0.1\,\times$ \code{resolution} of the surface, sharp edges are bevelled at the size of the
    grid, and concave corners are rounded with a radius of $0.75\,\times$ \code{resolution}.
  \item Parts thinner than about $1.5\,\times$ \code{resolution} vanish at \code{offset=0}. A
    loose surface has no thickness and needs an \code{offset} of at least $0.75\,\times$
    \code{resolution}. If nothing is left, \code{ValueError} is raised.
  \item \code{fill_holes=True} (default) fills everything that cannot be reached from outside,
    so only the outer hull remains. With \code{False}, enclosed voids stay as inner components
    with normals pointing into the void. This includes the inside of a closed shell, since the
    input is treated as a set of surfaces: the result is then a wall of thickness
    $2\,\times$ \code{offset}.
  \item The grid covers the bounding box. \code{max_nodes} (default 20 million) limits its
    size; the \code{ValueError} raised beyond it tells the size the grid would have had.
\end{itemize}

\paragraph{Method} The input is sampled into points a quarter of the resolution apart, and the
distance to them is computed on the grid in a narrow band. The nodes that can be reached from
the border of the grid without coming closer to the input than the hull distance are outside.
The surface between them and the rest is extracted with marching tetrahedra, which always gives
a closed, consistently oriented 2-manifold. That hull is watertight only at a distance of at
least $0.75\,\times$ \code{resolution}; for a smaller \code{offset}, a second distance field
measured from the hull gives the surface that lies the difference further in (a morphological
closing).

\subsection{Triangle quality and remeshing}

The triangles cut out of the grid can be arbitrarily thin. \apiref{mesh.SurfaceMesh.quality}
gives the shape quality of every triangle, $4\sqrt{3}\,A / (a^2 + b^2 + c^2)$ with the area $A$
and the sides $a$, $b$, $c$: 1 for an equilateral triangle and 0 for a degenerate one.
\code{shrink_wrap} checks it. If the worst triangle is below \code{min_quality} (default 0.3),
the nodes are moved within the surface towards the centres of their neighbours until it is
reached, in at most 50 smoothing steps. The connectivity is kept, so the mesh stays closed, and
since the nodes move in their tangent planes, the shape and the volume change by fractions of a
per cent of the resolution only. Typically 5 steps give a worst quality above 0.3 and a mean of
0.85. \code{min_quality=0} returns the mesh as cut out of the grid.

For triangles of a different size, or of still better shape (mean quality above 0.95),
\apiref{mesh.remesh} builds a new mesh through gmsh:

\begin{lstlisting}[style=python]
from nurbforge.mesh import remesh

coarse = remesh(hull, size=0.5)        # angle=40: edges sharper than 40 degrees are kept
\end{lstlisting}

gmsh splits the mesh into patches along its sharp edges, parametrises them and meshes them
anew; the new nodes lie on the old triangles. The result is again a watertight,
outward-oriented \code{SurfaceMesh}. \code{remesh} refuses meshes with nearly degenerate
triangles (quality below 0.01, \code{ValueError}), on which gmsh's classification does not
terminate.

\section{Convenience functions}

\paragraph{\apiref{mesh.mesh_surfaces}\code{(surfaces, path, size=None, order=1)}} Adds
surfaces and trimmed surfaces, meshes them in 2D and writes the result: a one-liner for quick
surface meshes of loose surfaces, such as lens surfaces.

\begin{lstlisting}[style=python]
from nurbforge.mesh import mesh_surfaces
mesh_surfaces([front, back, edge], "lens.msh", size=1.0)
\end{lstlisting}

\paragraph{\apiref{mesh.export_step}\code{(objects, path, solid=True)}} Writes surfaces,
trimmed surfaces and shells through OCC as STEP, BREP or IGES (by extension). Closed shells
become solids if \code{solid} is set. See \cref{sec:step-gmsh} for the comparison with the
native writer.

\paragraph{\apiref{mesh.bem_mesh}} See \cref{sec:mesh-bem}.

\section{Output formats}\label{sec:mesh-output}

\code{g.write(path)} calls \code{gmsh.write}, which chooses the format from the extension. The
most important formats are:

\begin{center}
\begin{tabularx}{\textwidth}{@{}lX@{}}
\toprule
Extension & Content \\
\midrule
\code{.msh} & gmsh's native mesh format, including physical groups. It is read by most
  open-source solvers (FEniCS/DOLFINx via \code{meshio} or \code{gmshio}, Elmer,
  bempp-cl, \dots). \\
\code{.vtk} & Legacy VTK, for ParaView and other visualisation tools. \\
\code{.stl} & Triangles only. The orientation is preserved; groups are lost. \\
\code{.inp}, \code{.unv}, \code{.bdf}, \code{.mesh}, \dots & Abaqus, I-DEAS, Nastran, Medit
  and others; see the gmsh documentation. \\
\code{.step}, \code{.brep}, \code{.iges} & The \emph{geometry} of the model (OCC) instead of
  the mesh. \\
\bottomrule
\end{tabularx}
\end{center}

Many gmsh options (\code{Mesh.MshFileVersion}, \code{Mesh.SaveAll}, \code{Mesh.Binary},
\dots) can be set with \code{g.gmsh.option.setNumber(...)} before \code{write}.

\section{Implementation notes}\label{sec:mesh-impl}

\paragraph{Shells travel through STEP.} \code{add_shell} writes the shell with \NF{}'s native
STEP writer (\cref{ch:step}) into a temporary file. It then reads the file with OCC's STEP
reader (\code{occ.importShapes}), with \code{Geometry.OCCTargetUnit} temporarily set to
\code{MM} so that no unit conversion takes place. The transfer stays exact, because the file
contains the original B-splines. OCC builds the topology itself, including the seam edges of
closed surfaces.

Why not transfer face by face? A trimmed closed face is a problem: a cylinder with a bore or a
slot, or a sphere, cone or torus cut by a Boolean. Its outer trim loop runs along the seam
twice. If the faces are added one by one and then merged with \code{removeAllDuplicates},
gmsh deletes such faces, and sewing leaves broken volumes. The STEP route avoids this.

\paragraph{Faces are found by interior points.} After the import, each face of the shell is
matched to the gmsh surface closest to an interior point of the face. The same matching is
used by \code{face_tags}, so labels remain valid after any renumbering.

\paragraph{\texttt{mesh(heal=True)}} merges duplicate points and edges of separately added
patches. This is harmless for shells, which are already consistent.

\paragraph{OCC's integrals.} OCC's \code{getMass} is only accurate to about 0.5\,\% for
trimmed rational faces. \NF{}'s own \code{area()} and \code{volume()} are accurate to the
requested tolerance. The volume of a fine tetrahedral mesh of a re-imported solid agrees with
\NF{}'s value, not with OCC's.

\paragraph{Labels and Booleans.} Labels are a kernel feature (\cref{ch:solids}), and
\code{nurbforge.mesh} only translates them into groups. A typical modelling workflow is
therefore:
\begin{enumerate}
  \item label the faces of the primitives;
  \item combine the primitives with Booleans;
  \item mesh with \code{physical_groups=True} and assign boundary conditions by name in the
    solver.
\end{enumerate}
